%! AP Statistics — Unit 2, Lesson 2.8 — Guided Notes
\documentclass[10pt]{article}
\usepackage{apstats-article}
\usepackage{apstats-boxes}

\begin{document}

\pageheader{Unit 2, Lesson 2.8}{Guided Notes}
\namedateperiod

\vspace{0.12in}

% ── OBJECTIVE ────────────────────────────────────────────────────────────────
\begin{objectivebox}
\small By the end of this lesson, I will be able to\dots\\[4pt]
\writeline\\[1pt]
\writeline\\[1pt]
\writeline
\end{objectivebox}

\vspace{0.1in}

% ── VOCABULARY ───────────────────────────────────────────────────────────────
\begin{vocabbox}
\termblanklong{Least-squares criterion}
\termblanklong{Coefficient of determination ($r^2$)}
\termblanklong{Influential point}
\termblanklong{High leverage point}
\end{vocabbox}

\vspace{0.1in}

% ── HOOK ─────────────────────────────────────────────────────────────────────
\begin{hookbox}
\small\textbf{Opening Question:} Your teacher uses a movable line on a scatterplot and shows
the sum of squared residuals changing as the line moves. Where should the line be to make that
sum as small as possible?

\vspace{6pt}
Sketch your best-guess line: (imagine a scatterplot here) \quad The sum of squared residuals
is smallest when the line is \blank{6cm}.

\vspace{4pt}
\textbf{Key idea:} There is an exact formula for this optimal line --- the
\textbf{least-squares regression line}.
\end{hookbox}

\vspace{0.1in}

% ── STEP 1 ───────────────────────────────────────────────────────────────────
\begin{notesbox}{Step 1 --- The Least-Squares Criterion (DAT-1.D)}
\small

\noindent The \textbf{least-squares criterion}: choose the line that minimizes

\[
  \sum e_i^2 \;=\; \sum \bigl(y_i - \hat{y}_i\bigr)^2
  \;=\; \blank{7cm}
\]

\vspace{4pt}
Why squaring? Because the sum of residuals $\sum e_i$ is always \blank{1.5cm} for any line
through $(\bar{x}, \bar{y})$, so squaring \blank{5.5cm} and penalizes
\blank{4cm} more than small ones.

\vspace{8pt}
\textbf{Result:} There is exactly \blank{2.5cm} least-squares line for any dataset.

\end{notesbox}

\vspace{0.1in}

% ── STEP 2 ───────────────────────────────────────────────────────────────────
\begin{notesbox}{Step 2 --- Computing Slope and Intercept (DAT-1.D)}
\small

\noindent Given summary statistics $\bar{x}$, $\bar{y}$, $s_x$, $s_y$, and $r$:

\vspace{6pt}
\textbf{Slope:} \quad $b = r\,\dfrac{s_y}{s_x}$ \quad\quad (\blank{3.5cm} of $r$ always matches \blank{2.5cm} of slope)

\vspace{8pt}
\textbf{Intercept:} \quad $a = \bar{y} - b\bar{x}$ \quad\quad (because the LSRL passes through \blank{3cm})

\vspace{10pt}
\textbf{Class example} --- hours of sleep per night ($x$) vs.\ reaction time in ms ($\hat{y}$):

\vspace{4pt}
$\bar{x} = 7.2$ hr, \; $\bar{y} = 310$ ms, \; $s_x = 1.5$ hr, \; $s_y = 45$ ms, \; $r = -0.80$

\vspace{6pt}
Step 1 --- Compute slope: \quad $b = \blank{5cm} = \blank{3cm}$

\vspace{6pt}
Step 2 --- Compute intercept: \quad $a = \blank{6cm} = \blank{3cm}$

\vspace{6pt}
Step 3 --- Write equation using variable names:

\vspace{4pt}
\hspace{1.5cm} \blank{10cm}

\end{notesbox}

\vspace{0.1in}

% ── STEP 3 ───────────────────────────────────────────────────────────────────
\begin{notesbox}{Step 3 --- The Coefficient of Determination $r^2$ (DAT-1.G)}
\small

\noindent $r^2$ is the \blank{3.5cm} of total variation in $y$ that is \blank{4.5cm}
by the linear model.

\vspace{6pt}
\textbf{AP template:} ``About \blank{2cm}\% of the variation in
\blank{3cm} is explained by the least-squares regression of
\blank{3cm} on \blank{3cm}.''

\vspace{8pt}
\textbf{Range:} $0 \leq r^2 \leq 1$.
\quad $r^2 = 1$: \blank{5cm}. \quad $r^2 = 0$: \blank{5cm}.

\vspace{8pt}
\textbf{For the sleep example:} $r = -0.80$, so $r^2 = \blank{2.5cm}$.

\vspace{4pt}
Interpret in context:\\[6pt]
\writeline\\[2pt]\writeline

\vspace{6pt}
\begin{tcolorbox}[colback=redbg, colframe=redacc, boxrule=0.8pt, arc=2mm,
    left=3mm, right=3mm, top=2mm, bottom=2mm]
\small\textbf{AP error to avoid:} ``The data are 64\% linear'' is \textbf{wrong.}
The correct phrasing always says ``variation in [response]'' and ``explained by.''
\end{tcolorbox}

\end{notesbox}

\vspace{0.1in}

% ── STEP 4 ───────────────────────────────────────────────────────────────────
\begin{notesbox}{Step 4 --- Influential Points (DAT-1.H)}
\small

\noindent An \textbf{influential point} is a point whose \blank{4cm} would
substantially \blank{4.5cm} the regression line.

\vspace{6pt}
A \textbf{high leverage point} has an extreme \blank{1.5cm} value (far from $\bar{x}$).
High leverage \blank{5cm} (circle one: always / does not always) create influence.

\vspace{8pt}
\textbf{How influential points affect the model:}

\vspace{4pt}
\begin{tabular}{@{} p{3.5cm} p{9cm} @{}}
Slope $b$ & \writeline \\[6pt]
Intercept $a$ & \writeline \\[6pt]
$r^2$ & \writeline \\[6pt]
\end{tabular}

\vspace{8pt}
\textbf{Rule:} An influential point should be \blank{3cm} (deleted / investigated) and
reported, not automatically removed.

\end{notesbox}

\vspace{0.1in}

% ── CONNECTIONS ───────────────────────────────────────────────────────────────
\begin{spiralbox}
\small
\begin{multicols}{2}

\textbf{How this connects:}
\begin{itemize}[itemsep=3pt]
  \item Lesson 2.5: $r$ drives the slope formula $b = r\frac{s_y}{s_x}$.
  \item Lesson 2.6: $a$ and $b$ give the regression equation $\hat{y} = a + bx$.
  \item Lesson 2.7: residuals are what least squares minimizes.
  \item Unit 9: inference for the slope uses the same LSRL framework.
\end{itemize}

\columnbreak

\textbf{Reflection:}\\
Why would an extreme $x$ value (high leverage) be more likely to be influential than an
extreme $y$ value?\\[2pt]
\writeline\\[1pt]\writeline\\[1pt]\writeline

\end{multicols}
\end{spiralbox}

\vspace{0.1in}

\begin{notesbox}{Additional Notes}
\vspace{0pt}
\writeline\\\writeline\\\writeline\\\writeline\\\writeline\\\writeline
\end{notesbox}

\end{document}
